\documentclass[../../main.tex]{subfiles}

\graphicspath{{../../image/}{../../figure/}}

\begin{document}

\section{习题}\label{section:ch15-习题}

\begin{exercise}\label{exercise:ch15-01}
（$\star$）考虑\cref{section:ch15-K均值聚类}讨论的 K 均值算法。证明由于离散指示变量 $r_{nk}$ 存在有限数量的可能取值，对于每个这样的取值都有一个唯一的最优集合 $\{\bm{\mu}_k\}$，因此 K 均值算法在有限次迭代后必定会收敛。
\end{exercise}

\begin{exercise}\label{exercise:ch15-02}
（$\star\star$）推导 K 均值算法的顺序形式。在每一步，考虑一个新的数据点 $\bm{x}_n$，并且只更新与数据点 $\bm{x}_n$ 最近的原型向量。从批处理设置中用于原型向量的\eqref{align:ch15-kmeans-mu}出发，将来自最后一个数据点 $\bm{x}_n$ 的贡献单独分离出来。通过重新排列各项，证明这种更新采取的是\eqref{align:ch15-kmeans-online}的形式。注意，由于在此推导过程中没有做任何近似，得到的原型向量具有以下性质：它们各自等于分配给它们的所有数据向量的均值。
\end{exercise}

\begin{exercise}\label{exercise:ch15-03}
（$\star$）考虑一个 Gauss 混合模型，其中潜变量的边缘分布 $p(\bm{z})$ 由\eqref{align:ch15-prior-factorized}给出，观测变量的条件分布 $p(\bm{x}\,|\,\bm{z})$ 由\eqref{align:ch15-conditional-x}给出。证明边缘分布 $p(\bm{x})$（通过对所有可能的 $\bm{z}$ 值求和 $p(\bm{z})p(\bm{x}\,|\,\bm{z})$ 而获得），是形式为\eqref{align:ch15-gmm-marginal}的 Gauss 混合分布。
\end{exercise}

\begin{exercise}\label{exercise:ch15-04}
（$\star$）证明在具有 $K$ 个分量的混合模型中，交换对称性导致的等价参数设置的数量是 $K!$。
\end{exercise}

\begin{exercise}\label{exercise:ch15-05}
（$\star\star$）假设我们希望使用 EM 算法来最大化包含潜变量的模型的参数后验分布 $p(\theta\,|\,\bm{X})$，其中 $\bm{X}$ 是观测数据集。证明 E 步骤在最大似然情况下保持不变，而 M 步骤中要最大化的量由 $\mathcal{Q}(\theta,\theta^{\text{old}})+\ln p(\theta)$ 给出，其中 $\mathcal{Q}(\theta,\theta^{\text{old}})$ 由\eqref{align:ch15-q-function}定义。
\end{exercise}

\begin{exercise}\label{exercise:ch15-06}
（$\star$）考虑\cref{figure:ch15-gmm-plate}所示的 Gauss 混合模型的有向图。利用 d 分离准则，证明潜变量的后验分布可以关于不同的数据点进行因子分解，即
\begin{align}
p(\bm{Z}\,|\,\bm{X},\bm{\mu},\bm{\Sigma},\bm{\pi})=\prod_{n=1}^{N}p(\bm{z}_n\,|\,\bm{x}_n,\bm{\mu},\bm{\Sigma},\bm{\pi}).
\label{align:ch15-ex-factor}
\end{align}
\end{exercise}

\begin{exercise}\label{exercise:ch15-07}
（$\star\star$）考虑 Gauss 混合模型的一种特殊情况，其中每个分量的协方差矩阵 $\bm{\Sigma}_k$ 都受限于它们所具有的一个共同的值 $\bm{\Sigma}$。在这样的模型下推导用于最大化似然函数的 EM 方程。
\end{exercise}

\begin{exercise}\label{exercise:ch15-08}
（$\star\star$）验证对 Gauss 混合模型的完整数据对数似然函数［\eqref{align:ch15-complete-nll}］进行最大化给出的 EM 方程。
\end{exercise}

\begin{exercise}\label{exercise:ch15-09}
（$\star\star$）证明如果我们在关于 $\bm{\Sigma}_k$ 最大化\eqref{align:ch15-q-gmm}的同时保持 $\gamma(z_{nk})$ 固定，我们将获得\eqref{align:ch15-mu-update}给出的闭式解。
\end{exercise}
\begin{remark}（重构新增：原文问题说明）
本题印为“关于 $\bm{\Sigma}_k$ 最大化”，却断言得到\eqref{align:ch15-mu-update}——后者是\emph{均值}的更新式。二者必有一处笔误：关于 $\bm{\Sigma}_k$ 求极值应得到协方差更新\eqref{align:ch15-sigma-update}（这恰是下一题 15.10 的内容），而得到\eqref{align:ch15-mu-update}的应是关于 $\bm{\mu}_k$ 求极值。这里保持原文照译，并在解答中把两种读法都完整推导一遍：按“关于 $\bm{\mu}_k$”读，结论\eqref{align:ch15-mu-update}成立；按“关于 $\bm{\Sigma}_k$”读，得到的是\eqref{align:ch15-sigma-update}。读者应以“关于 $\bm{\mu}_k$”为准。
\end{remark}

\begin{exercise}\label{exercise:ch15-10}
（$\star\star$）证明如果我们在关于 $\bm{\Sigma}_k$ 和 $\pi_k$ 最大化\eqref{align:ch15-q-gmm}的同时保持 $\gamma(z_{nk})$ 固定，我们将获得\eqref{align:ch15-sigma-update}和\eqref{align:ch15-pi-update}给出的闭式解。
\end{exercise}

\begin{exercise}\label{exercise:ch15-11}
（$\star\star$）考虑由混合分布给出的概率密度模型：
\begin{align}
p(\bm{x})=\sum_{k=1}^{K}\pi_k p(\bm{x}\,|\,k)
\label{align:ch15-ex-mixture}
\end{align}
假设我们将向量 $\bm{x}$ 分割成两部分，即 $\bm{x}=(\bm{x}_a,\bm{x}_b)$。

证明条件密度 $p(\bm{x}_a\,|\,\bm{x}_b)$ 本身就是一个混合分布，并给出混合系数和分量密度的表达式。
\end{exercise}

\begin{exercise}\label{exercise:ch15-12}
（$\star$）在\cref{subsection:ch15-em与kmeans}小节中，我们通过考虑一个所有分量都有协方差 $\varepsilon\bm{I}$ 的混合模型，得到了 K 均值算法和 EM 算法之间的关系。证明在 $\varepsilon\rightarrow0$ 的情况下，最大化这个模型的完整数据对数似然函数的期望［\eqref{align:ch15-q-gmm}］等同于最小化 K 均值算法给出的误差度量 $J$［\eqref{align:ch15-kmeans-loss}］。
\end{exercise}

\begin{exercise}\label{exercise:ch15-13}
（$\star\star$）验证 Bernoulli 分布的均值和协方差结果［\eqref{align:ch15-bernoulli-mean}和\eqref{align:ch15-bernoulli-cov}］。
\end{exercise}

\begin{exercise}\label{exercise:ch15-14}
（$\star\star$）考虑如下形式的混合分布：
\begin{align}
p(\bm{x})=\sum_{k=1}^{K}\pi_k p(\bm{x}\,|\,k)
\label{align:ch15-ex-mixture2}
\end{align}
其中 $\bm{x}$ 的元素可以是离散变量或连续变量，也可以是它们两者的组合。$p(\bm{x}\,|\,k)$ 的均值和协方差分别为 $\bm{\mu}_k$ 和 $\bm{\Sigma}_k$。利用\cref{exercise:ch15-13}的结果，证明该混合分布的均值和协方差分别由\eqref{align:ch15-bmm-mean}和\eqref{align:ch15-bmm-cov}给出。
\end{exercise}

\begin{exercise}\label{exercise:ch15-15}
（$\star\star$）使用 EM 算法的重新估计方程，证明当其参数设置对应于似然函数的最大值时，Bernoulli 混合模型具有以下性质：
\begin{align}
\mathbb{E}[\bm{x}]=\frac{1}{N}\sum_{n=1}^{N}\bm{x}_n=\bar{\bm{x}}.
\label{align:ch15-ex-mean}
\end{align}
继续证明如果这个模型的参数初始化使得所有分量都有相同的平均值 $\bm{\mu}_k=\hat{\bm{\mu}}$，$k=1,\cdots,K$，则 EM 算法将在任何初始混合系数的选择下，在一次迭代后收敛，且这个解具有 $\bm{\mu}_k=\bar{\bm{x}}$ 的性质。注意，这代表了混合模型的一种退化情况，实践中我们可以尝试通过进行适当的初始化来避免这种解。
\end{exercise}

\begin{exercise}\label{exercise:ch15-16}
（$\star$）考虑通过 $p(\bm{x}\,|\,\bm{z},\bm{\mu})$［\eqref{align:ch15-bmm-conditional}］和 $p(\bm{z}\,|\,\bm{\pi})$［\eqref{align:ch15-bmm-prior}］的乘积而得到的 Bernoulli 分布的潜变量和观测变量的联合分布。证明如果我们就 $\bm{z}$ 对这个联合分布进行边缘化，我们将得到\eqref{align:ch15-bernoulli-mixture}。
\end{exercise}

\begin{exercise}\label{exercise:ch15-17}
（$\star$）证明如果针对 $\bm{\mu}_k$ 最大化 Bernoulli 混合模型的完整数据对数似然函数的期望［\eqref{align:ch15-bmm-q}］，我们将得到 M 步骤方程［\eqref{align:ch15-bmm-mu}］。
\end{exercise}

\begin{exercise}\label{exercise:ch15-18}
（$\star$）证明如果我们针对混合系数 $\pi_k$ 最大化 Bernoulli 混合模型的完整数据对数似然函数的期望［\eqref{align:ch15-bmm-q}］，并使用 Lagrange 乘子引入求和约束，我们将得到 M 步骤方程［\eqref{align:ch15-bmm-pi}］。
\end{exercise}

\begin{exercise}\label{exercise:ch15-19}
（$\star$）证明由于离散变量 $\bm{x}_n$ 的概率 $p(\bm{x}_n\,|\,\bm{\mu}_k)$ 满足约束条件 $0\leqslant p(\bm{x}_n\,|\,\bm{\mu}_k)\leqslant1$，Bernoulli 混合模型的不完整数据对数似然函数存在上界，因此不存在使对数似然函数值趋于无穷大的奇异点。
\end{exercise}

\begin{exercise}\label{exercise:ch15-20}
（$\star\star\star$）考虑一个 $D$ 维变量 $\bm{x}$，它的每个分量本身是一个 $M$ 阶多项式变量，因此 $\bm{x}$ 是一个二值向量，其分量为 $x_{ij}$，其中 $i=1,\cdots,D$ 且 $j=1,\cdots,M$，约束条件 $\sum_jx_{ij}=1$ 对所有 $i$ 都成立。假设这些变量的分布可以用一系列离散多项式分布的混合分布来描述（见\cref{subsection:ch03-多项分布}），则有
\begin{align}
p(\bm{x})=\sum_{k=1}^{K}\pi_k\,p(\bm{x}\,|\,\bm{\mu}_k),
\label{align:ch15-ex-polymix}
\end{align}
其中
\begin{align}
p(\bm{x}\,|\,\bm{\mu}_k)=\prod_{i=1}^{D}\prod_{j=1}^{M}\mu_{kij}^{x_{ij}}.
\label{align:ch15-ex-polycomp}
\end{align}
参数 $\mu_{kij}$ 表示概率 $p(x_{ij}=1\,|\,\bm{\mu}_k)$，并且必须满足 $0\leqslant\mu_{kij}\leqslant1$，约束条件 $\sum_j\mu_{kij}=1$ 对所有的 $k$ 值和 $i$ 值都成立。

给定一个观测数据集 $\{\bm{x}_n\}$，其中 $n=1,\cdots,N$，推导出 EM 算法中的 E 步骤方程和 M 步骤方程，以便通过最大似然法优化混合系数 $\pi_k$ 和分量参数 $\mu_{kij}$。
\end{exercise}

\begin{exercise}\label{exercise:ch15-21}
（$\star$）验证\eqref{align:ch15-elbo-decomp}中的 $\mathcal{L}(q,\theta)$ 和 $\operatorname{KL}(q\,\|\,p)$ 分别由\eqref{align:ch15-elbo}和\eqref{align:ch15-kl}定义。
\end{exercise}

\begin{exercise}\label{exercise:ch15-22}
（$\star$）证明\eqref{align:ch15-elbo}给出的下界 $\mathcal{L}(q,\theta)$ 以 $q(\bm{Z})=p(\bm{Z}\,|\,\bm{X},\theta^{\text{old}})$ 的形式存在，其关于 $\theta$ 的梯度与对数似然函数 $\ln p(\bm{X}\,|\,\theta)$ 在 $\theta=\theta^{\text{old}}$ 处的梯度相同。
\end{exercise}

\begin{exercise}\label{exercise:ch15-23}
（$\star\star$）考虑应用于 Gauss 混合模型的 EM 算法的增量形式，在此形式中，只针对特定数据点 $\bm{x}_m$ 重新计算责任。从 M 步公式［\eqref{align:ch15-mu-update}和\eqref{align:ch15-effective-n}］出发，推导更新分量均值的结果［\eqref{align:ch15-online-mu}和\eqref{align:ch15-online-n}］。
\end{exercise}

\begin{exercise}\label{exercise:ch15-24}
（$\star\star$）当责任值以增量方式更新时，推导出在 Gauss 混合模型中更新协方差矩阵和混合系数的 M 步骤公式，这类似于更新分量均值的结果［\eqref{align:ch15-online-mu}］。
\end{exercise}

\end{document}
